\labsection{3}{Reshape time-aware data and engineer defensible features}{sec:lab03}

\begin{labproject}{Wrangling and feature engineering: build an analysis-ready table}
\textbf{Coverage.} Chapter~\ref{ch:wrangling}.\\
\textbf{Goal.} You should be able to preserve observational grain, construct derived variables, aggregate intentionally, validate engineered features, and explain what information is lost when data are summarised.\\
\textbf{Dataset.} \texttt{bike\_sharing\_dataset\_hour.csv}.\\
\textbf{Scenario.} Hourly bike records are too granular for a manager who wants to understand commuting demand. Your task is to transform the source table into a compact, traceable analytical dataset without losing the meaning of the original measurements.

\textbf{Core workflow.}
\begin{enumerate}
  \item \textbf{State the grain.} Confirm that one row represents one observed hour on one date. Verify the date and hour ranges before any transformation.
  \item \textbf{Preserve raw variables.} Create readable day-type and hour labels as new fields rather than replacing coded source columns.
  \item \textbf{Engineer an operational feature.} Define peak-period categories with explicit hour boundaries. Explain why the chosen boundaries are analytical assumptions rather than facts stored in the data.
  \item \textbf{Represent cyclical time.} Create sine/cosine encodings for hour-of-day so that 23:00 is close to 00:00 in feature space. Verify the unit-circle identity numerically.
  \item \textbf{Aggregate for a question.} Produce mean rental profiles by hour and working-day status. State how the observational unit changes after aggregation.
  \item \textbf{Separate demand components.} Summarise \texttt{casual} and \texttt{registered} rentals by peak period and day type. Identify where the two user groups show materially different patterns.
  \item \textbf{Validate the transformation.} Check that group counts sum to the source row count, engineered categories contain no unexpected missing values, and total rentals remain consistent with their components.
\end{enumerate}

\begin{tryit}
\textbf{Extension.} Change one peak-period boundary and quantify how much the resulting group summaries change. Use the result to explain why feature engineering should be treated as a modelling decision that may require sensitivity analysis.
\end{tryit}

\tcbbreak
\lstinputlisting[style=utpr]{all_experiments/lab03_wrangling.R}
\begin{verification}
\textbf{Required evidence.} Submit the final analysis-ready column definitions, one grouped demand profile, and three validation checks. For every engineered variable, include a one-sentence definition precise enough for another analyst to reproduce it independently.

\textbf{Decision checkpoint.} Recommend one feature that should enter a future demand model and one feature that should \emph{not} be used without additional care. Justify both choices from the variable meaning and the risk of leakage or redundancy.
\end{verification}
\end{labproject}
